\section{Arquitectura Multimodal en Tiempo Real}
\label{sec:arquitectura}

La extensión vive en un módulo aislado (\texttt{realtime\_lab/}) construido sobre una
ventana nativa (OpenCV \cite{bradski2000opencv} $+$ \texttt{sounddevice}) con tres
hilos concurrentes y colas seguras entre ellos.
El principio rector es un \textbf{puente de solo lectura}: los modelos del experimento
base (agentes, TME, encoder, decoder, estadísticas de cuantización) se cargan una sola
vez al inicio y nunca se escriben; el único estado que aprende en vivo es un
\emph{directorio de sesión} que existe solo en memoria.

\subsection{Componentes}

\begin{figure}[H]
\centering
\begin{tikzpicture}[
    node distance=6mm and 10mm,
    caja/.style={draw, rounded corners=2pt, align=center, font=\small,
                 minimum height=8mm, inner sep=4pt},
    hilo/.style={caja, fill=gray!12},
    mem/.style={caja, fill=gray!30},
    flecha/.style={-{Stealth[length=2.5mm]}, thick}
]
\node[hilo] (audio) {Hilo de audio\\ \footnotesize 16 kHz mono, bloques de 30 ms};
\node[caja, below=of audio] (vad) {VAD por energía + Whisper \emph{small}\\
    \footnotesize subtítulo parcial $\sim$1/s;\\
    \footnotesize pausa $\sim$600 ms $\to$ transcripción (\texttt{translate})};
\node[hilo, right=22mm of audio] (vision) {Hilo de visión\\
    \footnotesize 1 análisis cada $\sim$250 ms};
\node[caja, below=of vision] (enc) {recorte central ajustable\\
    \footnotesize $\to 128^2 \to$ ResNet18 $\to$ latente cuantizado};
\node[mem, below=10mm of $(vad.south)!0.5!(enc.south)$] (router)
    {\textbf{Router (solo lectura)}\\
     \footnotesize texto: broadcast TME (temprano) o $M_{\mathrm{dir}}$ del agente
     de entrada, lectura B1 (maduro)\\
     \footnotesize visión: $M^{R}_{\mathrm{dir}}$ del agente de entrada,
     lectura tolerante $\xi=2$\\
     \footnotesize evocación asíncrona: SS $+$ prototipo $+$ distancias};
\node[hilo, below=8mm of router] (ui) {Hilo de interfaz ($\sim$26--30 fps)\\
    \footnotesize video $+$ etiqueta viva $+$ subtítulos $+$ panel de sistema
    (redirecciones, barras, recuerdo evocado, $M_{\mathrm{dir}}$ de sesión)};
\draw[flecha] (audio) -- (vad);
\draw[flecha] (vision) -- (enc);
\draw[flecha] (vad) -- (router);
\draw[flecha] (enc) -- (router);
\draw[flecha] (router) -- (ui);
\end{tikzpicture}
\caption{Arquitectura de tres hilos. Las traducciones de modalidad (voz$\to$texto,
pixeles$\to$latente) ocurren antes del router; dentro del router solo hay operaciones
de memoria.}
\label{fig:arquitectura}
\end{figure}

\subsection{Canal de voz: los dos modos transactivos}

Cada frase transcrita se tokeniza y vectoriza con el pipeline oficial del experimento
(spaCy \cite{honnibal2020spacy} $+$ fastText con cuantización por magnitud) y se
enruta según el modo activo, conmutable en vivo:

\begin{description}
    \item[Modo temprano (TME activo).] Difusión a los ocho agentes: cada uno puntúa la
    pista con su score gateado por containment (el promedio de activación de la
    proyección, cero si alguna fila queda sin soporte) y el grupo asigna al máximo.
    Cada pista se \emph{registra} a nombre del ganador en el directorio de sesión, cuyo
    conteo por agente es visible al pie del panel: el directorio se forma frente al
    usuario, frase a frase.

    \item[Modo maduro (TME apagado).] Protocolo punto a punto del experimento base: el
    \emph{agente de entrada} (seleccionable en vivo) consulta su directorio textual
    entrenado con la lectura normalizada B1 y redirige al especialista ---o rechaza si
    su directorio no tiene señal (\texttt{directory\_no\_support}).
    El panel muestra la relación \textsc{entrada} $\to$ \textsc{destino} con la
    anotación ``redirige'' o ``se queda (es el experto)''.
\end{description}

El rechazo es de primera clase en ambos modos: una frase sin tokens representables o
sin soporte en ninguna memoria produce \textsc{rechazo} explícito con su causa, nunca
un ganador por desempate.

\subsection{Canal de visión: etiqueta viva y percepción honesta}

El cuadro de la cámara pasa por un recorte cuadrado central \emph{ajustable}
(55\% del lado menor por defecto), se reescala a $128\times128$ y se codifica con el
encoder ResNet18 del experimento; el latente cuantizado consulta el directorio visual
$M^{R}_{\mathrm{dir}}$ del agente de entrada con la lectura tolerante ($\xi = 2$) del
protocolo oficial.
La etiqueta sobre el video se suaviza por mayoría sobre los últimos 5 análisis para
evitar parpadeo, y el rechazo muestra siempre el candidato más cercano con su score.

Dos decisiones de interfaz resultaron determinantes en las pruebas físicas
(Sección~\ref{sec:met_fisica}):

\begin{itemize}
    \item \textbf{El recuadro ajustable.} La primera versión exigía llenar todo el
    alto del video ($480$ px), lo que obligaba a acercar el estímulo impreso más allá
    de la distancia de enfoque de la cámara web. El recuadro al 55\% se llena a
    20--30 cm, donde la óptica sí enfoca.
    \item \textbf{``Lo que ve el ojo''.} Un recuadro en vivo con el $128\times128$
    exacto que recibe el encoder. Desenfoque, reflejo o desencuadre dejan de ser
    hipótesis: se ven.
\end{itemize}

\subsubsection*{Política de color: reintentar con lentes}

La cadena física impresora$+$cámara desplaza los momentos de color (imágenes grisáceas,
brillo corrido) y el latente resultante cae fuera del soporte del directorio.
La mitigación es una renormalización por canal hacia los momentos de referencia de
ETH-80 (media y desviación RGB promediadas sobre imágenes de entrenamiento):
\begin{equation}
    x'_c \;=\; \frac{x_c - \mu_c}{\sigma_c + \epsilon}\,\hat{\sigma}_c + \hat{\mu}_c,
    \qquad c \in \{R, G, B\}.
    \label{eq:colorfix}
\end{equation}
La política es deliberadamente \emph{sin umbral}: intento 1 con el cuadro crudo; solo
si el directorio rechaza, intento 2 con la corrección~\eqref{eq:colorfix}.
Se descartó un detector previo por distancia de momentos porque, medido, las
distribuciones de entradas limpias y degradadas se traslapan (mediana 24.2 vs.\ 27.1):
no existe el umbral que las separe.
El reintento condicionado no puede dañar una entrada limpia ---esta ya fue aceptada en
el intento crudo--- y la interfaz marca \texttt{[color-fixed]} cuando el rescate actuó.

\subsection{Evocación asíncrona: prototipo y objeto definido}
\label{sec:arq_pista}

Lo único lento del sistema es la recuperación estocástica (SS: 127 muestras más
descenso local, $\sim$2 s).
Para no congelar video ni voz, la evocación corre en un hilo aparte: el panel muestra
``recalling\ldots'' y la imagen aparece al estar lista.

La evocación implementa el contraste de la Sección~\ref{sec:pista_faltante}.
Para la pista ganadora se calculan:
(i) el \textbf{objeto definido}: el patrón devuelto por el recall oficial (SS),
decodificado a imagen;
(ii) el \textbf{prototipo emergente}~\eqref{eq:prototipo} del mismo plano proyectado,
decodificado a imagen; y
(iii) la \textbf{distancia de retro-proyección}~\eqref{eq:test} de ambos.
El panel los apila con sus distancias: la masa leída columna a columna arriba
(borde gris), el objeto definido abajo (borde del color del agente).
Solo se muestra la reconstrucción de memoria ---nunca la fotografía ETH-80 de
referencia al lado--- para que no parezca que la imagen proviene del encoder en lugar
de la memoria.

\subsection{Los tres métodos en la aplicación de trazado}

En paralelo a la ventana en tiempo real, la aplicación de trazado del experimento
(la interfaz Streamlit del pipeline, etapa de recall) incorpora la suite completa:
para el token que produce la imagen final se calculan RS, ST y SS sobre la
\emph{misma} proyección, junto con el prototipo, y se muestran las cuatro imágenes con
sus distancias y el plano proyectado como mapa de calor
($64$ características $\times$ $32$ valores).
Los métodos se implementaron en la envoltura del sistema
(\texttt{recall\_from\_left($\cdot$, method)}) reutilizando las primitivas de la
implementación de referencia (\texttt{project}, \texttt{reduce},
\texttt{distance\_recall}); el método por defecto sigue siendo SS, de modo que ningún
resultado previo del sistema cambia.
